% ============================================================================
% SEC04.TEX - Section 3.4: Kontrol Alur Program
% ============================================================================

\section{Kontrol Alur Program}

\begin{learningoutcomes}
\item Menggunakan if-else untuk kondisi
\item Menggunakan switch-case
\item Menggunakan loop (for, while, foreach)
\item Menerapkan kontrol alur dalam konteks game
\end{learningoutcomes}

\subsection{If-Else Statement}

\begin{lstlisting}[language={[Sharp]C}]
if (health > 0)
{
    // Player masih hidup
}
else
{
    // Player mati
    GameOver();
}
\end{lstlisting}

\subsection{Switch-Case}

\begin{lstlisting}[language={[Sharp]C}]
switch (gameState)
{
    case GameState.Menu:
        ShowMenu();
        break;
    case GameState.Playing:
        StartGame();
        break;
    case GameState.Paused:
        PauseGame();
        break;
    default:
        Debug.Log("Unknown state");
        break;
}
\end{lstlisting}

\subsection{Loop}

\begin{lstlisting}[language={[Sharp]C}]
// For loop
for (int i = 0; i < 10; i++)
{
    Debug.Log(i);
}

// While loop
while (health > 0)
{
    // Game continues
}

// Foreach loop
foreach (GameObject enemy in enemies)
{
    enemy.SetActive(false);
}
\end{lstlisting}
